\section{Method}

The staged baseline performs deterministic agglomerative merging while tracking
merge witnesses. Each witness records enough information to explain why regions
were combined and to distinguish valid construction from infeasible capacity
cases.

The atlas reading rule is that larger regions are not self-explanatory. The
paper must show the adjacent-pair candidate set, the eligibility gate, the
score used only after eligibility is established, and the merge witness that
creates the next non-leaf region. Rejected adjacent pairs remain evidence when
they explain why an apparently plausible merge was not taken.

\subsection{Reproducible Procedure}

The constructor operates on the adjacency graph emitted by the BISECT state
pipeline. Connectedness is defined over that graph: candidate regions and final
districts must be connected under the declared profile. Islands, disconnected
components, and alternative adjacency definitions are therefore input-profile
questions rather than claims solved by the merge routine.

\begin{enumerate}
  \item Initialize each unit as a singleton region with population and
        adjacency metadata.
  \item Enumerate adjacent candidate region pairs that remain eligible under
        the declared capacity and connectedness profile.
  \item Score candidate merges by capacity eligibility first, then distance of
        merged population from the ideal district population.
  \item Select the best candidate, breaking ties by merged population and
        stable region identifiers.
  \item Record a merge witness and repeat until the requested number of
        regions is reached or no eligible merge remains.
\end{enumerate}

\begin{table}[h]
\centering
\small
\begin{tabular}{ll}
\toprule
Witness field & Purpose \\
\midrule
\texttt{step} & Orders the merge in the hierarchy \\
\texttt{left\_region}, \texttt{right\_region} & Identifies the adjacent regions selected for merge \\
\texttt{merged\_region} & Links the new region to its descendants \\
\texttt{merged\_population} & Shows the population created by the merge \\
\texttt{cut\_edges\_between} & Records adjacency strength between the merged regions \\
\bottomrule
\end{tabular}
\caption{Implemented merge-witness fields for replayable regionalization
lineage. The summary sidecar separately records merge count, hierarchy depth,
capacity status, repair status, edge cut, and parameter hash.}
\end{table}

\subsection{Worked Merge Decision}

The merge step is a local election among adjacent region pairs. Suppose the
ideal district population is \(100\) with an upper capacity of \(110\). The
constructor first separates over-capacity pairs from eligible pairs, then ranks
eligible pairs by closeness to \(100\). A geographically adjacent pair can still
lose if it would exceed capacity.

\begin{table}[h]
\centering
\small
\begin{tabular}{lllll}
\toprule
Pair & Merged pop. & Over capacity? & Distance to ideal & Decision \\
\midrule
A+B & 118 & yes & 18 & defer: over capacity \\
B+C & 96 & no & 4 & selected \\
C+D & 88 & no & 12 & eligible, lower rank \\
\bottomrule
\end{tabular}
\caption{T.16 chooses among adjacent candidate merges. Capacity eligibility is
checked before the score is used.}
\end{table}
